\paragraph{Prescribed-frequency quartet theorem.}
Let each branch of a four-taxon species tree carry a finite piecewise-constant
frequency history ordered from its younger to older boundary. Conditional on
these histories, assume background switching rates
\(\rho p\) and \(\rho(1-p)\), same-background coalescence rates
\(c_0(p)=1/[M\max(1-p,1/M)]\) and
\(c_1(p)=1/[M\max(p,1/M)]\), and the absorbing genealogy generators
\(G_n(p)\) used by the MSRC simulator. Then the conditional quartet law is
obtained by applying, on every branch, the ordered product
\[
\prod_j\exp(G_n(p_j)\Delta t_j)
\]
to the appropriate young-boundary row distribution. For the balanced tree,
\[
q_i=q_{\mathrm{early},i}+F_{\mathrm{bal}}^{\mathsf T}
P_4(\mathcal P_A)h_i,
\]
and for the pectinate tree,
\[
q_i=c_i+F_{\mathrm{pect}}^{\mathsf T}P_4(\mathcal P_A)h_i.
\]
Here (c_i) is the mass absorbed into topology (i) before the four-lineage
root stage. The formulas are exact conditional on the supplied histories and
conserve probability. If the histories are random Moran histories, the
unconditional law is the expectation of these conditional expressions; one
prescribed trajectory is not an unconditional finite-(M) result.

\paragraph{Proof.}
The transient and absorbing states form a finite continuous-time Markov chain.
The forward row Kolmogorov equation is (v'(t)=v(t)G_n(p)), giving the stated
matrix exponential on each constant segment. Concatenation of segments gives
the ordered product. At each speciation boundary, the conditional history
supplies the daughter-boundary background distributions, so the product of
independent daughter contributions is the conditional joint distribution.
For the balanced tree, failure of either sister pair to survive contributes
only (Q_1); the surviving mass is (F_{\rm bal}), after which the root
chain absorbs into (h_i). For the pectinate tree, the three-lineage chain
absorbs into (c_i) or enters the root as (F_{\rm pect}). These events are
disjoint and exhaustive, proving the formulas and probability conservation.
Taking expectation over the random histories proves the final statement.
